\documentclass[11pt]{article}
\usepackage[utf8]{inputenc}
\usepackage[T1]{fontenc}
\usepackage[brazil]{babel}
\usepackage{geometry}
\usepackage{tcolorbox}
\usepackage{ragged2e}

\geometry{a4paper, margin=1.8cm}
\tcbset{colback=white}

\begin{document}

\begin{center}
\Large \textbf{Revisão – Cultura Digital} \\
\large \textbf{Segurança Digital e Ética Online}
\\[0.5cm]

\textbf{Professor: Jefferson}
\end{center}

%%%%%%%%%%%%%%%%%%%%%%%%%%%%%%%%%%%%%%%%%%%%%%%%%%%%%%%%%%%%%%
%%%%%%%%%%%%%%%%%%%%%%%% QUESTÃO 1 %%%%%%%%%%%%%%%%%%%%%%%%%%%%
%%%%%%%%%%%%%%%%%%%%%%%%%%%%%%%%%%%%%%%%%%%%%%%%%%%%%%%%%%%%%%

\begin{tcolorbox}[colback=blue!10,colframe=blue!60!black,
title=\textbf{Questão 1 – Senha Forte}]
O que é uma senha forte e por que é importante utilizá-la em todas as suas contas digitais?
Dê um exemplo de uma senha segura.
\end{tcolorbox}

\begin{tcolorbox}[colback=green!10,colframe=green!60!black,
title=\textbf{Solução 1 — Explicação }]
Uma \textbf{senha forte} é aquela difícil de adivinhar, usando letras maiúsculas, minúsculas,
números e símbolos. Ela protege suas contas contra invasões e roubo de dados.

\textbf{Exemplo de senha forte:} \texttt{Fz!92kP\#11}
\end{tcolorbox}

%%%%%%%%%%%%%%%%%%%%%%%%%%%%%%%%%%%%%%%%%%%%%%%%%%%%%%%%%%%%%%
%%%%%%%%%%%%%%%%%%%%%%%% QUESTÃO 2 %%%%%%%%%%%%%%%%%%%%%%%%%%%%
%%%%%%%%%%%%%%%%%%%%%%%%%%%%%%%%%%%%%%%%%%%%%%%%%%%%%%%%%%%%%%

\begin{tcolorbox}[colback=blue!10,colframe=blue!60!black,
title=\textbf{Questão 2 – Phishing}]
O que é phishing e como podemos nos proteger desse golpe?  
Descreva uma situação em que isso poderia acontecer.
\end{tcolorbox}

\begin{tcolorbox}[colback=green!10,colframe=green!60!black,
title=\textbf{Solução 2 — Explicação }]
\textbf{Phishing} é um golpe em que criminosos fingem ser empresas ou pessoas confiáveis
para roubar dados, como senhas e informações bancárias.  
Para se proteger, é importante verificar links, desconfiar de mensagens urgentes
e nunca informar senhas por e-mail.

\textbf{Exemplo de situação:}  
Um e-mail falso dizendo que seu banco bloqueará sua conta se você não “confirmar seus dados”.
\end{tcolorbox}

%%%%%%%%%%%%%%%%%%%%%%%%%%%%%%%%%%%%%%%%%%%%%%%%%%%%%%%%%%%%%%
%%%%%%%%%%%%%%%%%%%%%%%% QUESTÃO 3 %%%%%%%%%%%%%%%%%%%%%%%%%%%%
%%%%%%%%%%%%%%%%%%%%%%%%%%%%%%%%%%%%%%%%%%%%%%%%%%%%%%%%%%%%%%

\begin{tcolorbox}[colback=blue!10,colframe=blue!60!black,
title=\textbf{Questão 3 – Privacidade nas Redes Sociais}]
Quais cuidados você deve ter ao compartilhar informações pessoais nas redes sociais?
Dê dois exemplos de boas práticas.
\end{tcolorbox}

\begin{tcolorbox}[colback=green!10,colframe=green!60!black,
title=\textbf{Solução 3 — Explicação }]
É importante evitar expor dados sensíveis e controlar quem tem acesso às suas publicações.

\textbf{Boas práticas:}
\begin{itemize}
\item Não compartilhar endereço, escola ou rotina diária.
\item Ajustar a privacidade das postagens.
\end{itemize}
\end{tcolorbox}

%%%%%%%%%%%%%%%%%%%%%%%%%%%%%%%%%%%%%%%%%%%%%%%%%%%%%%%%%%%%%%
%%%%%%%%%%%%%%%%%%%%%%%% QUESTÃO 4 %%%%%%%%%%%%%%%%%%%%%%%%%%%%
%%%%%%%%%%%%%%%%%%%%%%%%%%%%%%%%%%%%%%%%%%%%%%%%%%%%%%%%%%%%%%

\begin{tcolorbox}[colback=blue!10,colframe=blue!60!black,
title=\textbf{Questão 4 – Cyberbullying}]
Explique o que é cyberbullying e como ele pode afetar as pessoas.  
O que você faria se visse alguém sofrendo agressão online?
\end{tcolorbox}

\begin{tcolorbox}[colback=green!10,colframe=green!60!black,
title=\textbf{Solução 4 — Explicação }]
\textbf{Cyberbullying} é agressão, humilhação ou ofensa praticada pela internet.
Ele pode causar tristeza profunda, ansiedade, isolamento e baixa autoestima.

Diante dessa situação, a atitude correta é apoiar a vítima, denunciar o agressor
e avisar um responsável ou autoridade.
\end{tcolorbox}

%%%%%%%%%%%%%%%%%%%%%%%%%%%%%%%%%%%%%%%%%%%%%%%%%%%%%%%%%%%%%%
%%%%%%%%%%%%%%%%%%%%%%%% QUESTÃO 5 %%%%%%%%%%%%%%%%%%%%%%%%%%%%
%%%%%%%%%%%%%%%%%%%%%%%%%%%%%%%%%%%%%%%%%%%%%%%%%%%%%%%%%%%%%%

\begin{tcolorbox}[colback=blue!10,colframe=blue!60!black,
title=\textbf{Questão 5 – Fake News}]
O que são fake news e como identificar se uma notícia é falsa?
Apresente dois critérios de verificação.
\end{tcolorbox}

\begin{tcolorbox}[colback=green!10,colframe=green!60!black,
title=\textbf{Solução 5 — Explicação }]
\textbf{Fake news} são notícias falsas criadas para enganar, manipular ou gerar confusão.

\textbf{Critérios para identificar:}
\begin{itemize}
\item Verificar a fonte e se outros veículos confiáveis confirmam a informação.
\item Conferir data, autor e possíveis sinais de manipulação.
\end{itemize}
\end{tcolorbox}

%%%%%%%%%%%%%%%%%%%%%%%%%%%%%%%%%%%%%%%%%%%%%%%%%%%%%%%%%%%%%%
%%%%%%%%%%%%%%%%%%%%%%%% QUESTÃO 6 %%%%%%%%%%%%%%%%%%%%%%%%%%%%
%%%%%%%%%%%%%%%%%%%%%%%%%%%%%%%%%%%%%%%%%%%%%%%%%%%%%%%%%%%%%%

\begin{tcolorbox}[colback=blue!10,colframe=blue!60!black,
title=\textbf{Questão 6 – Antivírus Atualizado}]
Por que é importante manter o antivírus atualizado?
O que pode acontecer se isso não for feito?
\end{tcolorbox}

\begin{tcolorbox}[colback=green!10,colframe=green!60!black,
title=\textbf{Solução 6 — Explicação }]
O antivírus atualizado reconhece novas ameaças.  
Se não for atualizado, o dispositivo pode ser infectado por vírus recentes,
causando perda de arquivos, roubo de dados ou travamentos.
\end{tcolorbox}

%%%%%%%%%%%%%%%%%%%%%%%%%%%%%%%%%%%%%%%%%%%%%%%%%%%%%%%%%%%%%%
%%%%%%%%%%%%%%%%%%%%%%%% QUESTÃO 7 %%%%%%%%%%%%%%%%%%%%%%%%%%%%
%%%%%%%%%%%%%%%%%%%%%%%%%%%%%%%%%%%%%%%%%%%%%%%%%%%%%%%%%%%%%%

\begin{tcolorbox}[colback=blue!10,colframe=blue!60!black,
title=\textbf{Questão 7 – Backup}]
O que é um backup e por que ele é importante?  
Dê um exemplo de situação em que o backup seria útil.
\end{tcolorbox}

\begin{tcolorbox}[colback=green!10,colframe=green!60!black,
title=\textbf{Solução 7 — Explicação }]
\textbf{Backup} é uma cópia de segurança dos arquivos.  
Ele é importante porque permite recuperar dados após danos, invasões,
roubos ou falhas no dispositivo.

\textbf{Exemplo:} Perder o celular e recuperar fotos e documentos através da nuvem.
\end{tcolorbox}

%%%%%%%%%%%%%%%%%%%%%%%%%%%%%%%%%%%%%%%%%%%%%%%%%%%%%%%%%%%%%%
%%%%%%%%%%%%%%%%%%%%%%%% QUESTÃO 8 %%%%%%%%%%%%%%%%%%%%%%%%%%%%
%%%%%%%%%%%%%%%%%%%%%%%%%%%%%%%%%%%%%%%%%%%%%%%%%%%%%%%%%%%%%%

\begin{tcolorbox}[colback=blue!10,colframe=blue!60!black,
title=\textbf{Questão 8 – Mídias Digitais}]
As mídias digitais mudaram a comunicação.  
Explique duas mudanças positivas e uma negativa dessas tecnologias.
\end{tcolorbox}

\begin{tcolorbox}[colback=green!10,colframe=green!60!black,
title=\textbf{Solução 8 — Explicação }]
\textbf{Duas mudanças positivas:}
\begin{itemize}
\item Comunicação rápida com qualquer pessoa.
\item Acesso a informações e conteúdos educativos.
\end{itemize}

\textbf{Mudança negativa:}  
Exposição excessiva e dependência das redes sociais.
\end{tcolorbox}

%%%%%%%%%%%%%%%%%%%%%%%%%%%%%%%%%%%%%%%%%%%%%%%%%%%%%%%%%%%%%%
%%%%%%%%%%%%%%%%%%%%%%%% QUESTÃO 9 %%%%%%%%%%%%%%%%%%%%%%%%%%%%
%%%%%%%%%%%%%%%%%%%%%%%%%%%%%%%%%%%%%%%%%%%%%%%%%%%%%%%%%%%%%%

\begin{tcolorbox}[colback=blue!10,colframe=blue!60!black,
title=\textbf{Questão 9 – Respeito à Privacidade}]
O que significa respeitar a privacidade alheia nas redes sociais?  
Dê um exemplo de atitude respeitosa e outra de atitude inadequada.
\end{tcolorbox}

\begin{tcolorbox}[colback=green!10,colframe=green!60!black,
title=\textbf{Solução 9 — Explicação }]
Respeitar a privacidade significa não divulgar informações, imagens ou conversas
de outras pessoas sem permissão.

\textbf{Exemplo respeitosa:} Pedir autorização antes de postar foto com alguém.  
\textbf{Exemplo inadequada:} Compartilhar prints de conversas privadas.
\end{tcolorbox}

%%%%%%%%%%%%%%%%%%%%%%%%%%%%%%%%%%%%%%%%%%%%%%%%%%%%%%%%%%%%%%
%%%%%%%%%%%%%%%%%%%%%%%% QUESTÃO 10 %%%%%%%%%%%%%%%%%%%%%%%%%%%
%%%%%%%%%%%%%%%%%%%%%%%%%%%%%%%%%%%%%%%%%%%%%%%%%%%%%%%%%%%%%%

\begin{tcolorbox}[colback=blue!10,colframe=blue!60!black,
title=\textbf{Questão 10 – Uso Ético da Internet}]
O que significa ser um usuário responsável e ético nas mídias digitais?  
Dê sua explicação e um exemplo de comportamento adequado.
\end{tcolorbox}

\begin{tcolorbox}[colback=green!10,colframe=green!60!black,
title=\textbf{Solução 10 — Explicação }]
Ser um usuário ético é agir com respeito, sinceridade e responsabilidade online,
evitando ofensas, mentiras e conteúdos prejudiciais.

\textbf{Exemplo de comportamento adequado:}  
Verificar informações antes de repassar e tratar todos com respeito.
\end{tcolorbox}

\end{document}

